% Knapsack interdiction: monotonicity, count formulation, pruning rules,
% online relaxation, robust tables, threshold form, per-node rule, exactness.
\section{Knapsack interdiction proofs}\label{app:kip}

This section proves the rules behind the exact search for subset-sum
knapsack interdiction (KIP with $p=w$, Section~\ref{app:notation}). The search
answers a yes/no question: given an integer $\ell$, does every leader set
within budget leave the follower at least $\ell$? A ``yes'' proves
$\mathrm{OPT}\ge\ell$; a ``no'' comes with a leader set whose follower value is
below~$\ell$. Repeating with $\ell$ equal to the incumbent value proves
optimality (Corollary~\ref{cor:kip-loop}).

\subsection{Monotonicity and the count formulation}

\begin{lemma}[Monotonicity of interdiction]\label{lem:kip-monotone}
Let $\mathcal{F}$ be any family of subsets of $N$ (the follower's feasible
sets) and $\Phi_{\mathcal{F}}(X)=\max\{p(Y): Y\in\mathcal{F},\ Y\cap X=\emptyset\}$
with $p\ge 0$ (a maximum over no set is $0$). If $X\subseteq X'$ then $\Phi_{\mathcal{F}}(X')\le\Phi_{\mathcal{F}}(X)$. If
$\mathcal{F}$ is closed under taking subsets, then
$\Phi_{\mathcal{F}}(X)=\max\{p(Y\setminus X): Y\in\mathcal{F}\}$. The knapsack
follower $\mathcal{F}=\{Y: \sum_{j\in Y}W_{\cdot j}\le c\}$ is closed under taking
subsets because $W\ge 0$; in particular $\Phi(X')\le\Phi(X)$ for $X\subseteq X'$.
\end{lemma}
\begin{proof}
Every $Y\in\mathcal{F}$ that avoids $X'$ also avoids $X$, so the maximum for $X'$
is over a subfamily. For the second claim, a set $Y\in\mathcal{F}$ avoiding $X$
has $p(Y)=p(Y\setminus X)$, which gives ``$\le$''; conversely, for
$Y\in\mathcal{F}$ the set $Y\setminus X$ lies in $\mathcal{F}$ and avoids $X$, which
gives ``$\ge$''. A subset of a packing has no larger weight in any row
because $W\ge 0$.
\end{proof}
{\raggedright\noindent Formalised: \texttt{Bobi/Knapsack.lean}, theorems
\texttt{follower\_antitone}, \texttt{bestAvoiding\_antitone},
\texttt{bestAvoiding\_eq\_sup\_sdiff}, \texttt{packings\_downward\_closed} and
\texttt{K\_eq\_bestAvoiding}.\par}

\begin{lemma}[Cheapest copies first]\label{lem:kip-cheapest}
Consider KIP and a class $Q\subseteq N$ of identical items ($p_i=p_j$ and
$w_i=w_j$ for $i,j\in Q$). For every leader set $X$ there is a leader set $X'$
with $X'\setminus Q=X\setminus Q$, $|X'\cap Q|=|X\cap Q|$,
$\gamma_a\le\gamma_b$ for all $a\in X'\cap Q$ and $b\in Q\setminus X'$,
$\gamma(X')\le\gamma(X)$ and $\Phi(X')=\Phi(X)$.
\end{lemma}
\begin{proof}
First a single swap: if $i\in X\cap Q$, $j\in Q\setminus X$ and
$\gamma_j\le\gamma_i$, let $X^{s}=(X\setminus\{i\})\cup\{j\}$. The transposition of
$i$ and $j$ preserves $p$ and $w$ and maps $N\setminus X$ onto $N\setminus X^{s}$,
so it maps the packings avoiding $X$ bijectively onto those avoiding $X^{s}$
with the same profits; hence $\Phi(X^{s})=\Phi(X)$, and
$\gamma(X^{s})=\gamma(X)-\gamma_i+\gamma_j\le\gamma(X)$. Now induct on the integer
$\gamma(X)$. If some $i\in X\cap Q$ and $j\in Q\setminus X$ have $\gamma_j<\gamma_i$,
the swap lowers the cost strictly, keeps $X\setminus Q$, $|X\cap Q|$ and
$\Phi$, and the induction hypothesis applies to $X^{s}$. Otherwise $X'=X$
works.
\end{proof}
{\raggedright\noindent Formalised: \texttt{Bobi/Knapsack.lean}, theorems
\texttt{K\_map\_equiv}, \texttt{swap\_block} and \texttt{cheapest\_first}.\par}

\begin{definition}[Class instance]\label{def:class-instance}
A \emph{class instance} is a list of classes $1,\dots,L$; class $i$ has a
value $v_i\in\mathbf{Z}_{\ge 1}$ and $n_i$ copies with deletion costs
$\gamma_{i,1},\dots,\gamma_{i,n_i}\in\mathbf{Z}_{\ge 0}$ (in a fixed order, not
necessarily ascending). Let $P_i(r)=\sum_{h\le r}\gamma_{i,h}$ for
$0\le r\le n_i$. A \emph{deletion} is a count vector $r=(r_1,\dots,r_L)$ with
$0\le r_i\le n_i$; its cost is $P(r)=\sum_i P_i(r_i)$ and its survivors form
the multiset $R(r)$ with $n_i-r_i$ copies of $v_i$ for each $i$. For a finite
multiset $M$ of nonnegative integers, $\mathrm{SS}(M)$ is the set of sums of its
sub-multisets ($0\in\mathrm{SS}(M)$). Given integers $0\le lo\le hi$ and a
budget $\beta\ge 0$, the instance is \emph{safe} if every deletion with
$P(r)\le\beta$ has $\mathrm{SS}(R(r))\cap[lo,hi]\neq\emptyset$; a deletion within
budget with $\mathrm{SS}(R(r))\cap[lo,hi]=\emptyset$ is \emph{blocking}.
\end{definition}

\begin{proposition}[Reduction of subset-sum KIP]\label{prop:kip-reduction}
Consider subset-sum KIP with data $\gamma,\beta,w,c$ (all nonnegative
integers). Let $F=\{j:\gamma_j=0\}$, $\mathcal{L}=\{j\notin F: 1\le w_j\le c\}$ and
$g=\gcd\{w_j: j\in\mathcal{L}\}$ ($g=1$ if $\mathcal{L}=\emptyset$). Group
$\mathcal{L}$ into classes by the value $w_j/g$, with each class's costs sorted
ascending, and let $hi=\lfloor c/g\rfloor$. For a deletion $r$ let $X(r)$ be $F$
together with the $r_i$ cheapest items of each class. Then
\begin{enumerate}
\item $\gamma(X(r))=P(r)$ and $\Phi(X(r))=g\cdot\max\bigl(\mathrm{SS}(R(r))\cap[0,hi]\bigr)$;
\item for every leader set $X$ with $\gamma(X)\le\beta$ there is a deletion $r$
with $P(r)\le\gamma(X)$ and $\Phi(X(r))\le\Phi(X)$;
\item for every integer $\ell\ge 0$ with $lo=\lceil\ell/g\rceil\le hi$:
$\mathrm{OPT}\ge\ell$ if and only if the class instance is safe for the band
$[lo,hi]$ and budget $\beta$; and $r$ is blocking if and only if
$P(r)\le\beta$ and $\Phi(X(r))<\ell$.
\end{enumerate}
\end{proposition}
\begin{proof}
(1) Items of $F$ cost nothing. A packing avoiding $X(r)$ contains no item of
$F$; an item outside $F\cup\mathcal{L}$ has $w_j=0$ (then $p_j=0$) or $w_j>c$ (then
it is in no packing), so dropping such items changes neither weight nor
profit. Every weight in $\mathcal{L}$ is $g$ times its class value, so a set
$Y\subseteq\mathcal{L}$ has $w(Y)\le c$ iff its value sum is at most
$\lfloor c/g\rfloor=hi$, and the value sums of packings of $\mathcal{L}\setminus X(r)$
are exactly $\mathrm{SS}(R(r))$.
(2) Let $X_1=X\cup F$: $\gamma(X_1)=\gamma(X)$ and $\Phi(X_1)\le\Phi(X)$ by
Lemma~\ref{lem:kip-monotone}. Items outside $F\cup\mathcal{L}$ do not affect $\Phi$
(as in~(1)) and blocking them costs something, so remove them from $X_1$
(cost does not increase, $\Phi$ unchanged). Apply Lemma~\ref{lem:kip-cheapest}
to each class in turn. The result blocks in each class $i$ some $r_i$ copies,
each no costlier than every unblocked copy, so its cost in class $i$ is the
sum of the $r_i$ smallest costs, $P_i(r_i)$, and its follower value is that
of $X(r)$.
(3) Since $g s\ge\ell$ iff $s\ge lo$ for integers $s$, (1) gives:
$\Phi(X(r))\ge\ell$ iff $\mathrm{SS}(R(r))\cap[lo,hi]\neq\emptyset$. If the instance is
safe, (2) gives $\Phi(X)\ge\Phi(X(r))\ge\ell$ for every $X$ within budget. If a
deletion $r$ is blocking, $X(r)$ is within budget with $\Phi(X(r))<\ell$, so
$\mathrm{OPT}<\ell$.
\end{proof}
\noindent Not formalised as a whole; its two ingredients are
Lemmas~\ref{lem:kip-monotone} and~\ref{lem:kip-cheapest}.

If $\lceil\ell/g\rceil>hi$, the question is trivial: the leader set $F=X(0)$
costs nothing and, by part~(1), $\Phi(F)\le g\cdot hi<\ell$, so
$\mathrm{OPT}<\ell$ and the deletion $r=0$ is a witness. From now on we work with a class instance, a band
$[lo,hi]$ with $lo\le hi$ and a budget~$\beta$. All results below hold for
arbitrary nonnegative costs in the count formulation; where ascending costs
are needed (deletions of arbitrary copies) we say so.

\subsection{Windows of reachable sums and the stall test}

\begin{lemma}[Window extension]\label{lem:window-extend}
Let $M$ be a multiset with $[a,b]\subseteq\mathrm{SS}(M)$, $a\le b$. If a value $v$
satisfies $v\le b-a+1$ then $[a,b+v]\subseteq\mathrm{SS}(M\uplus\{v\})$. More
generally, if values $v_1,\dots,v_q$ satisfy
$v_h\le (b-a+1)+v_1+\dots+v_{h-1}$ for every $h$, then
$[a,\,b+\sum_h v_h]\subseteq\mathrm{SS}(M\uplus\{v_1,\dots,v_q\})$.
\end{lemma}
\begin{proof}
A target $t\in[a,b]$ is a sum of $M$. A target $t\in(b,b+v]$ has
$a\le t-v\le b$, because $t>b$ and $v\le b-a+1$; so $t$ is $v$ plus a sum of
$M$. The
general statement follows by induction on $q$: after adding $v_1,\dots,v_{h-1}$
the window is $[a,b+v_1+\dots+v_{h-1}]$, of width
$(b-a+1)+v_1+\dots+v_{h-1}\ge v_h$.
\end{proof}
{\raggedright\noindent Formalised: \texttt{Bobi/SubsetSum.lean}, theorems
\texttt{window\_extend} and \texttt{window\_extend\_all}.\par}

\begin{lemma}[Survivors that fit]\label{lem:window-survivors}
Let $[a,b]\subseteq\mathrm{SS}(M)$ with $a\le b$ and $\theta=b-a+1$. Let
$(v_1,\eta_1),\dots,(v_q,\eta_q)$ be a list of copies with values $v_h$ and
flags $\eta_h\in\{\text{kept},\text{deleted}\}$. Write $\mu_h=\sum_{h'<h}v_{h'}$
for the mass before copy $h$ and $\delta_h$ for the deleted mass before copy
$h$. If $\delta_h+v_h\le\theta+\mu_h$ for every $h$, then
$[a,\,b+\sigma]\subseteq\mathrm{SS}(M\uplus K)$, where $K$ is the multiset of kept
copies and $\sigma$ its total.
\end{lemma}
\begin{proof}
Let $h$ be a kept copy. The kept mass before it is $\mu_h-\delta_h$, and the
hypothesis says $v_h\le\theta+(\mu_h-\delta_h)$. So the kept copies, in list
order, satisfy the hypothesis of Lemma~\ref{lem:window-extend}.
\end{proof}
{\raggedright\noindent Formalised: \texttt{Bobi/SubsetSum.lean}, theorems
\texttt{fits\_survivors} and \texttt{window\_survivors}
(\texttt{SubsetSum.perm}: subset sums do not depend on the order).\par}

\begin{lemma}[Fractional deletion bound]\label{lem:fractional}
Let copies $h$ have values $v_h>0$ and costs $\gamma_h\ge 0$, sorted so that
$\gamma_1/v_1\le\gamma_2/v_2\le\cdots$. For $d>0$ let
$\Psi(d)=\min\{\sum_h\gamma_h z_h:\sum_h v_hz_h\ge d,\ 0\le z\le 1\}$ ($+\infty$
if $\sum_h v_h<d$). The greedy solution, which fills copies in this order
until mass $d$ is reached, is optimal, and every set $H$ of copies with
$\sum_{h\in H}v_h\ge d$ has $\sum_{h\in H}\gamma_h\ge\Psi(d)$.
\end{lemma}
\begin{proof}
The indicator vector of $H$ is feasible, which proves the last claim. Let
$m_h=v_hz_h$ and $\rho_h=\gamma_h/v_h$, so the cost is $\sum_h\rho_hm_h$ with
$0\le m_h\le v_h$ and $\sum_h m_h\ge d$. Let $m^{*}$ be the greedy allocation,
$\lambda$ the ratio of the copy where it stops, and $m$ any feasible
allocation. Then
$\sum_h\rho_h(m_h-m^{*}_h)=\sum_h(\rho_h-\lambda)(m_h-m^{*}_h)+\lambda\sum_h(m_h-m^{*}_h)$.
Each term of the first sum is nonnegative ($m^{*}_h=v_h\ge m_h$ when
$\rho_h<\lambda$, $m^{*}_h=0\le m_h$ when $\rho_h>\lambda$), and the last term is
nonnegative because $\sum_hm^{*}_h=d\le\sum_hm_h$ and $\lambda\ge 0$.
\end{proof}
\noindent Not formalised (a standard linear programming fact).

The next rule is stated for a search node (Section~\ref{app:kip-exact}): the
decided classes have left a multiset $M$ of survivors, $S=\mathrm{SS}(M)\cap[0,hi]$,
and the remaining (undecided) classes may be deleted with residual budget
$\beta'$. A \emph{completion} is a count vector for the undecided classes of
cost at most~$\beta'$.

\begin{proposition}[Window rule with stall test]\label{prop:window-rule}
Order the undecided classes by ascending value, $u_1,\dots,u_q$ (the proof
works for any order; ascending value is the code's choice), and let
$M_i=\sum_{i'<i}n_{u_{i'}}v_{u_{i'}}$ be the mass of the classes before $u_i$.
Let $\Delta$ be the largest mass that a completion deletes and $\sigma_0$ the
total undecided mass minus $\Delta$. Let $[a,b]\subseteq S$ be a run of
consecutive integers with $b<lo\le b+\sigma_0$ and width $\theta=b-a+1$, and
put $D_i=M_i-v_{u_i}+\theta+1$. Suppose that for every $i$: $D_i>0$, and either
$i=1$ or the fractional bound of Lemma~\ref{lem:fractional} over the copies
of $u_1,\dots,u_{i-1}$ satisfies $\Psi_i(D_i)>\beta'$. Then every completion
leaves survivors $R$ with $lo\in\mathrm{SS}(M\uplus R)$; the node is safe.
\end{proposition}
\begin{proof}
Fix a completion and list the undecided copies class by class in the order
$u_1,\dots,u_q$, flagging the deleted ones. Let $e_i$ be the mass the completion
deletes from $u_1,\dots,u_{i-1}$. Those deleted copies form a set of total cost
at most $\beta'$, so $e_i\ge D_i$ would contradict $\Psi_i(D_i)>\beta'$ by
Lemma~\ref{lem:fractional}; for $i=1$, $e_1=0<D_1$. Hence $e_i\le D_i-1$. At
the copy number $t$ ($t=0,1,\dots$) of class $u_i$, the mass before it is
$M_i+tv_{u_i}$ and the deleted mass before it is at most $e_i+tv_{u_i}$, so
deleted mass plus $v_{u_i}$ is at most
$D_i-1+(t+1)v_{u_i}=\theta+M_i+tv_{u_i}$. Lemma~\ref{lem:window-survivors}
gives $[a,b+\sigma]\subseteq\mathrm{SS}(M\uplus R)$ with $\sigma\ge\sigma_0$ the
surviving mass, and $a\le b<lo\le b+\sigma_0$ puts $lo$ in that window.
\end{proof}
{\raggedright\noindent The per-copy hypothesis is formalised
(\texttt{Bobi/SubsetSum.lean}, theorem \texttt{window\_survivors}); the
reduction to one test per class and the use of Lemma~\ref{lem:fractional} are
argued here and in the doc comment of the code, not formalised. $\Delta$ is
computed exactly by a bounded-knapsack recursion over the undecided classes.\par}

\subsection{The online relaxation}

The online relaxation of Weninger and Fukasawa\cite{weninger2023fast,weninger2025fast}
(Theorem~2 in Section~3 of their preprint) applies to general KIP. Fix a list
$o=(j_1,\dots,j_L)$ of distinct items, any subset of $N$ in any order. The
\emph{online game} on $o$ has $2L$ rounds: at item $j_h$ the leader first
decides whether to block it (allowed if its cost fits the remaining budget);
if the leader does not block it, the follower decides whether to take it
(allowed if its weight fits the remaining capacity). The follower collects
the profits of the items it takes. Both players see all earlier moves. The
value of this finite game of perfect information is given by backward
induction: $\Omega(\emptyset,\beta',\kappa)=0$ and, for $o=(j,o')$,
\begin{align*}
  \Omega(o,\beta',\kappa)&=
  \begin{cases}\min\{\Omega(o',\beta'-\gamma_j,\kappa),\ T\} & \text{if }\gamma_j\le\beta',\\
  T & \text{otherwise,}\end{cases}\\
  T&=\begin{cases}\max\{\Omega(o',\beta',\kappa),\ \Omega(o',\beta',\kappa-w_j)+p_j\} & \text{if }w_j\le\kappa,\\
  \Omega(o',\beta',\kappa) & \text{otherwise.}\end{cases}
\end{align*}

\begin{theorem}[Online lower bound, any play order]\label{thm:online}
For every list $o$ of distinct items and every $X\subseteq N$ with
$\gamma(X)\le\beta$,
$\Omega(o,\beta,c)\le K(o\setminus X)\le\Phi(X)$, where $K(T)$ is the best
packing of the items of $T$ with capacity $c$. Hence $\Omega(o,\beta,c)\le\mathrm{OPT}$.
\end{theorem}
\begin{proof}
In the online game the leader may adapt to the follower's earlier moves; the
real game corresponds to the leader strategy ``block exactly the items of
$X$'', which does not adapt. Formally, we prove by induction on the length of
$o$, for all $\beta',\kappa\ge 0$ and all $X\subseteq o$ with $\gamma(X)\le\beta'$,
that $\Omega(o,\beta',\kappa)\le K_\kappa(o\setminus X)$ (best packing with
capacity $\kappa$). For $o=\emptyset$ both sides are $0$. Let $o=(j,o')$. If
$j\in X$, then $\gamma_j\le\gamma(X)\le\beta'$, so the blocking branch is offered and
$\Omega(o,\beta',\kappa)\le\Omega(o',\beta'-\gamma_j,\kappa)\le K_\kappa(o'\setminus X)=K_\kappa(o\setminus X)$
by induction applied to $X\setminus\{j\}$. If $j\notin X$, then $X\subseteq o'$ and
$\Omega(o,\beta',\kappa)\le T$. By induction $\Omega(o',\beta',\kappa)\le K_\kappa(o'\setminus X)$
and, if $w_j\le\kappa$,
$\Omega(o',\beta',\kappa-w_j)+p_j\le K_{\kappa-w_j}(o'\setminus X)+p_j\le K_\kappa((o'\setminus X)\cup\{j\})$
(add $j$ to an optimal packing). So in both cases of $T$ we get
$T\le K_\kappa((o'\setminus X)\cup\{j\})=K_\kappa(o\setminus X)$. For an arbitrary $X$ within
budget apply this to $X\cap o$ (its cost is no larger and $o\setminus(X\cap o)=o\setminus X$).
Finally $K(o\setminus X)\le K(N\setminus X)=\Phi(X)$ because more items can only help.
\end{proof}
{\raggedright\noindent Formalised: \texttt{Bobi/Online.lean}, theorems
\texttt{online\_le\_K}, \texttt{online\_le\_K\_any} and
\texttt{online\_lower\_bound}.\par}

\begin{remark}
The bound need not be tight. In Weninger and Fukasawa's Table~1 (profits
$(1,1,2)$, leader costs $(1,1,2)$, follower weights $(1,1,3)$, $\beta=2$, $c=3$)
the online game in the play order $(1,2,3)$ is worth $1$ and the bilevel optimum is $2$ (Lean:
\texttt{table1\_online}, \texttt{table1\_opt}, \texttt{table1\_block}, checked by
evaluation).
\end{remark}

\subsection{Robust tables with an end set}

For subset-sum instances the same idea is applied class by class, in the
count formulation. Let $Q=(q_1,\dots,q_L)$ be a list of classes (the
\emph{play order}) and $A\subseteq\mathbf{Z}_{\ge 0}$ an \emph{end set}. The game
$\Gamma_A(Q,\beta',t)$ starts with residual $t$ and budget $\beta'$; at class $q$
the leader announces $r\le n_q$ with $P_q(r)$ at most the remaining budget,
then the follower takes $j\le n_q-r$ copies with $jv_q$ at most the current
residual, which decreases by $jv_q$; the follower wins if the final residual
lies in~$A$.

\begin{definition}[Robust tables]\label{def:tables}
$G_A(\emptyset,\beta')=A$ and, for $Q=(q,Q')$,
\[
  t\in G_A(Q,\beta')\iff \forall r\le n_q \text{ with } P_q(r)\le\beta'\ \
  \exists j\le n_q-r:\ jv_q\le t\ \text{and}\ t-jv_q\in G_A(Q',\beta'-P_q(r)).
\]
So $G_A(Q,\beta')$ is the set of residuals from which the follower wins
$\Gamma_A(Q,\beta',\cdot)$. The robust table is $H=G_{\{0\}}$ and the band table of
width $w\ge 1$ is $G_w=G_{[0,w-1]}$.
\end{definition}

\begin{theorem}[Soundness of the tables]\label{thm:tables-sound}
If $t\in G_A(Q,\beta')$ and $r$ is a deletion of the classes of $Q$ with
$P(r)\le\beta'$, then there are $s\in\mathrm{SS}(R(r))$ and $a\in A$ with $s+a=t$. In
particular $t\in H(Q,\beta')$ gives $t\in\mathrm{SS}(R(r))$, and $t\in G_w(Q,\beta')$ gives
some $s\in\mathrm{SS}(R(r))$ with $s\le t<s+w$. If each class's costs are ascending,
the same holds for a deletion of arbitrary copies of total cost at most
$\beta'$.
\end{theorem}
\begin{proof}
Induction on $L$. For $L=0$ take $s=0$, $a=t$. For $Q=(q,Q')$:
$P_q(r_1)\le P(r)\le\beta'$, so the definition gives $j\le n_q-r_1$ with
$t-jv_q\in G_A(Q',\beta'-P_q(r_1))$, and the rest of $r$ costs at most
$\beta'-P_q(r_1)$. By induction $t-jv_q=s'+a$ with $s'$ a subset sum of the
survivors of $Q'$; then $s=jv_q+s'$ uses $j$ of the $n_q-r_1$ survivors of
$q$. For deletions of arbitrary copies with ascending costs, deleting $r_i$
copies of class $i$ costs at least $P_i(r_i)$, and the survivors are again
$n_i-r_i$ copies of $v_i$.
\end{proof}
{\raggedright\noindent Formalised: \texttt{Bobi/EndTable.lean}, theorems
\texttt{GA\_sound}, \texttt{GA\_sound\_any}, \texttt{G\_iff\_GA},
\texttt{H\_iff\_GA}, \texttt{G\_sound\_of\_GA}, \texttt{H\_sound\_of\_GA};
\texttt{Bobi/RobustTable.lean}, \texttt{H\_sound}, \texttt{H\_sound\_any},
\texttt{take\_sum\_le}, \texttt{counts\_of\_deletion};
\texttt{Bobi/BandTable.lean}, \texttt{G\_sound}, \texttt{G\_sound\_any},
\texttt{G\_one}, \texttt{G\_sound\_one}.\par}

\begin{proposition}[Any play order]\label{prop:play-order}
Let $Q$ be the undecided classes in the search order and $\pi(Q)$ any
permutation of them. If $t\in G_A(\pi(Q),\beta')$, the conclusion of
Theorem~\ref{thm:tables-sound} holds for every deletion $r$ of $Q$ with
$P(r)\le\beta'$.
\end{proposition}
\begin{proof}
Permuting the pairs (class, count) permutes the survivors and leaves the cost
unchanged, and subset sums depend only on the multiset. Apply
Theorem~\ref{thm:tables-sound} to the permuted deletion.
\end{proof}
{\raggedright\noindent Formalised: \texttt{Bobi/EndTable.lean}, theorems
\texttt{counts\_perm}, \texttt{remaining\_eq\_flatten},
\texttt{totalCost\_eq\_sum}, \texttt{remaining\_perm},
\texttt{totalCost\_perm}, \texttt{GA\_sound\_perm} and
\texttt{GA\_sound\_perm\_any}.\par}

\begin{corollary}[Root-table rules]\label{cor:root-rules}
Let $M$ be the decided survivors and $Q$ the undecided classes (any play
order). (i) If $s\in\mathrm{SS}(M)$ and $t\in H(Q,\beta')$, every completion leaves
$s+t\in\mathrm{SS}(M\uplus R)$. (ii) If $[a,b]\subseteq\mathrm{SS}(M)$, $t\in G_w(Q,\beta')$,
$a+t\le\tau$ and $\tau+w\le b+t+1$, every completion leaves
$\tau\in\mathrm{SS}(M\uplus R)$. The search prunes when $s+t$, respectively $\tau$,
lies in $[lo,hi]$.
\end{corollary}
\begin{proof}
(i) Theorem~\ref{thm:tables-sound} gives $t\in\mathrm{SS}(R)$. (ii) It gives
$s'\in\mathrm{SS}(R)$ with $s'\le t<s'+w$; then $a\le\tau-t\le\tau-s'$ and
$\tau-s'\le\tau-t+w-1\le b$, so $\tau-s'\in\mathrm{SS}(M)$.
\end{proof}
{\raggedright\noindent Formalised: \texttt{Bobi/RobustTable.lean}, theorem
\texttt{covers\_sound}; \texttt{Bobi/BandTable.lean},
\texttt{band\_covers\_sound}, \texttt{covers\_sound\_of\_band};
\texttt{Bobi/Threshold.lean}, \texttt{covers\_threshold\_sound},
\texttt{covers\_threshold\_sound\_perm}, \texttt{band\_covers\_threshold\_sound},
\texttt{band\_covers\_threshold\_sound\_perm} (the same on stored thresholds).\par}

\begin{lemma}[Antitone in the budget]\label{lem:antitone}
If $\beta''\le\beta'$ then $G_A(Q,\beta')\subseteq G_A(Q,\beta'')$. If
$A\subseteq A'$ then $G_A(Q,\beta')\subseteq G_{A'}(Q,\beta')$.
\end{lemma}
\begin{proof}
Induction on $L$; the base case does not depend on the budget. A deletion
$r$ with $P_q(r)\le\beta''$ also has $P_q(r)\le\beta'$; the answer $j$ for
budget $\beta'$ leads to $G_A(Q',\beta'-P_q(r))\subseteq G_A(Q',\beta''-P_q(r))$ by
induction. The second claim is the same induction with a larger base set.
\end{proof}
{\raggedright\noindent Formalised: \texttt{Bobi/EndTable.lean}, theorems
\texttt{GA\_antitone}, \texttt{G\_antitone}, \texttt{H\_antitone} and
\texttt{GA\_mono}.\par}

\subsection{Threshold form}

By Lemma~\ref{lem:antitone}, for each residual $t$ the budgets at which
$t\in G_A(Q,\cdot)$ form an initial segment, so one number per $t$ describes
the table. Fix a cap $\mathrm{cap}\ge 0$ (the root budget $\beta$).

\begin{definition}[Threshold]\label{def:thr}
$\mathrm{thr}(\emptyset,t)=\mathrm{cap}$ if $t\in A$ and $-1$ otherwise; for
$Q=(q,Q')$ with value $v$ and $n$ copies,
\[
  \mathrm{thr}(Q,t)=\min\Bigl(\mathrm{cap},\ \min_{\substack{r\le n\\ P_q(r)\le\mathrm{cap}}}
  \Bigl[P_q(r)+\max\Bigl(-1,\ \max_{\substack{j\le n-r\\ jv\le t}}\mathrm{thr}(Q',t-jv)\Bigr)\Bigr]\Bigr).
\]
\end{definition}

\begin{theorem}[Threshold form]\label{thm:thr}
For every integer $B\ge 0$: $B\le\mathrm{thr}(Q,t)$ if and only if $B\le\mathrm{cap}$
and $t\in G_A(Q,B)$. Consequently $-1\le\mathrm{thr}(Q,t)\le\mathrm{cap}$;
$\mathrm{thr}(Q,t)=-1$ iff $t\notin G_A(Q,0)$; and otherwise $\mathrm{thr}(Q,t)$ is the
largest $B\le\mathrm{cap}$ with $t\in G_A(Q,B)$.
\end{theorem}
\begin{proof}
Induction on $L$. For $L=0$: if $t\in A$ the claim reads $B\le\mathrm{cap}$; if
$t\notin A$ both sides are false since $B\ge 0>-1$. For $Q=(q,Q')$ write
$m_r(t)$ for the inner maximum (at least $-1$). Then $B\le\mathrm{thr}(Q,t)$ iff
$B\le\mathrm{cap}$ and $B\le P_q(r)+m_r(t)$ for all $r$ with $P_q(r)\le\mathrm{cap}$.
Assume $B\le\mathrm{cap}$. For $r$ with $P_q(r)>B$ the inequality holds
automatically, as $B\le P_q(r)-1$. For $r$ with $P_q(r)\le B$ (these have
$P_q(r)\le\mathrm{cap}$): $B\le P_q(r)+m_r(t)$ iff $0\le B-P_q(r)\le m_r(t)$ iff some
$j\le n-r$ with $jv\le t$ has $B-P_q(r)\le\mathrm{thr}(Q',t-jv)$ iff, by induction
(note $B-P_q(r)\le\mathrm{cap}$), $t-jv\in G_A(Q',B-P_q(r))$. This is the
definition of $t\in G_A(Q,B)$. The consequences follow by taking $B=0$ and
$B=\mathrm{thr}(Q,t)$.
\end{proof}
{\raggedright\noindent Formalised: \texttt{Bobi/Threshold.lean}, theorems
\texttt{thr\_iff}, \texttt{GA\_iff\_thr}, \texttt{thr\_bounds},
\texttt{thr\_eq\_neg\_one\_iff}, \texttt{thr\_greatest}, \texttt{G\_iff\_thr}
and \texttt{H\_iff\_thr} (definition \texttt{thr}).\par}

\begin{remark}[What the code computes]
The recursion at residual $t$ reads only residuals $\le t$, so arrays indexed
$0,\dots,hi$ hold the exact thresholds of those residuals. With nonnegative
costs $P_q$ is nondecreasing, so stopping at the first $r$ with
$P_q(r)>\mathrm{cap}$ implements the filter. That the code's loops compute
$\mathrm{thr}$ (including the shortcuts that compute only residuals in a range
that can matter, and an early exit once every computed threshold is below the
budget) is argued in the code's doc comments and tested against a direct
transcription of Definition~\ref{def:thr} and against brute force; it is not
formalised.
\end{remark}

\subsection{The per-node rule}

At a search node the decided survivors $M$ reach the sums $S=\mathrm{SS}(M)\cap[0,hi]$.
The follower wins with any total in $[lo,hi]$, so it need not fix its target
in advance. The node's end set is
\[
  A_{\mathrm{node}}=\{a\in[0,hi] : \exists s\in S,\ s\le a\le s+(hi-lo)\},
\]
that is, $(S+[0,hi-lo])\cap[0,hi]$, and the node table asks whether
$hi\in G_{A_{\mathrm{node}}}(Q,\beta')$.

\begin{theorem}[Per-node rule]\label{thm:node-rule}
Let $S\subseteq\mathrm{SS}(M)$, $lo\le hi$, $A$ any set whose elements $a$ each have
some $s\in S$ with $s\le a\le s+(hi-lo)$, and $\pi(Q)$ any permutation of the
undecided classes $Q$. If $hi\in G_A(\pi(Q),\beta')$, then every completion
(a deletion $r$ of $Q$ with $P(r)\le\beta'$, costs in any order) leaves some
$u\in\mathrm{SS}(M\uplus R(r))$ with $lo\le u\le hi$. In threshold form: if
$\beta'\le\mathrm{cap}$ and $\beta'\le\mathrm{thr}(\pi(Q),hi)$ computed with end set
$A_{\mathrm{node}}$, the node is safe at every budget up to $\mathrm{thr}(\pi(Q),hi)$.
\end{theorem}
\begin{proof}
Proposition~\ref{prop:play-order} gives $x\in\mathrm{SS}(R(r))$ and $a\in A$ with
$x+a=hi$. Pick $s\in S$ with $s\le a\le s+(hi-lo)$. Then $u=s+x$ is a subset sum
of $M\uplus R(r)$ and $u=s+hi-a\in[hi-(hi-lo),hi]=[lo,hi]$. The threshold form
follows from Theorem~\ref{thm:thr} with $B=\beta'$ (and for smaller budgets,
since $B\le\mathrm{thr}$ is the same condition). The hypothesis $lo\le hi$ is
needed: with $lo>hi$ there is no target, and the code returns $-1$.
\end{proof}
{\raggedright\noindent Formalised: \texttt{Bobi/EndTable.lean}, theorems
\texttt{node\_rule}, \texttt{node\_table\_sound}, \texttt{node\_table\_sound\_perm}
and \texttt{node\_table\_sound\_any} (definition \texttt{nodeEnd});
\texttt{Bobi/Threshold.lean}, \texttt{node\_table\_threshold\_sound} and
\texttt{node\_table\_threshold\_sound\_counts} (the form the search uses:
count formulation, any play order, any cost order).\par}

\begin{remark}[Strength]
The per-node rule contains the rule of $H$ in the same play order: if $s\in S$,
$s\le\tau$, $\tau\in[lo,hi]$ and $\tau-s\in H(Q,\beta')$, then
$hi\in G_{A_{\mathrm{node}}}(Q,\beta')$ (shift the follower's strategy up by
$hi-\tau+s$). Formalised: \texttt{Bobi/EndTable.lean}, theorem
\texttt{node\_table\_of\_H}. The analogous containment of the band rules is
argued in the code's doc comment, not formalised; soundness does not need
either. The play order matters only for strength; the search plays the
classes in ascending order of deletion cost per unit of mass.
\end{remark}

\subsection{Restricted copies}

\begin{proposition}[Restricted follower]\label{prop:restrict}
For $k\ge 1$ let $I^{(k)}$ be the class instance in which every class $i$ with
$n_i>k$ is replaced by a class with the same value, $k$ copies and costs
$(\gamma_{i,1}+\dots+\gamma_{i,n_i-k+1},\ \gamma_{i,n_i-k+2},\dots,\gamma_{i,n_i})$. If
$I^{(k)}$ is safe (same band and budget), then so is the original instance.
\end{proposition}
\begin{proof}
Map a deletion $r$ of the original to $r'$ with
$r'_i=\max(0,r_i-(n_i-k))$ for the restricted classes and $r'_i=r_i$ otherwise.
If $r'_i=0$ its cost is $0\le P_i(r_i)$; otherwise it is
$\gamma_{i,1}+\dots+\gamma_{i,n_i-k+r'_i}=P_i(r_i)$. So $P'(r')\le P(r)$. The
restricted survivors number $k-r'_i=\min(k,n_i-r_i)\le n_i-r_i$, so
$R'(r')\subseteq R(r)$ as multisets and $\mathrm{SS}(R'(r'))\subseteq\mathrm{SS}(R(r))$. Safety
of $I^{(k)}$ applied to $r'$ gives a sum in $[lo,hi]$ for $r$.
\end{proof}
\noindent Not formalised. The restricted costs are not ascending, which is
why the count-formulation results above make no assumption on the order of
costs. (The option is off by default; no run recorded in the repository
uses it.)

\subsection{Why the search is exact}\label{app:kip-exact}

The search fixes an order of the classes (any order). A node at depth $k$ has
decided counts $r_1,\dots,r_k$ with residual budget $\beta'=\beta-\sum_{i\le k}P_i(r_i)\ge 0$,
decided survivors $M_k$ and $S_k=\mathrm{SS}(M_k)\cap[0,hi]$. The node is \emph{safe}
if every completion $(r_{k+1},\dots,r_L)$ with $\sum_{i>k}P_i(r_i)\le\beta'$ has
$\mathrm{SS}(M_k\uplus R_{>k})\cap[lo,hi]\neq\emptyset$. The root ($k=0$, $S_0=\{0\}$) is safe
iff the class instance is safe.

\begin{lemma}[Node state]\label{lem:node-state}
Whether a node is safe depends only on $(k,S_k,\beta')$, and a node safe at
$\beta'$ is safe at every $\beta''\le\beta'$.
\end{lemma}
\begin{proof}
$\mathrm{SS}(M\uplus R)=\mathrm{SS}(M)+\mathrm{SS}(R)$, and all values are nonnegative, so a sum in
$[lo,hi]$ has its $\mathrm{SS}(M)$ part in $[0,hi]$. Hence the condition reads
$(S_k+\mathrm{SS}(R_{>k}))\cap[lo,hi]\neq\emptyset$, and the completions depend only on $k$
and $\beta'$, with fewer completions for smaller budgets.
\end{proof}

\begin{theorem}[Sound pruning]\label{thm:kip-rules}
Each of the following implies that the node $(k,S_k,\beta')$ is safe.
\begin{enumerate}
\item Direct: $S_k\cap[lo,hi]\neq\emptyset$.
\item Guaranteed copies: with $d_i=\max\{r\le n_i: P_i(r)\le\beta'\}$, the
multiset $U$ with $n_i-d_i$ copies of $v_i$ for each undecided $i$ has
$(S_k+\mathrm{SS}(U))\cap[lo,hi]\neq\emptyset$.
\item Window: the hypothesis of Proposition~\ref{prop:window-rule} holds for a
run of $S_k$.
\item Root tables: Corollary~\ref{cor:root-rules} gives a target in $[lo,hi]$,
for a table built once at the root (any play order).
\item Node table: $\beta'\le\mathrm{thr}(\pi(Q),hi)$ with end set $A_{\mathrm{node}}$
(Theorem~\ref{thm:node-rule}).
\item Cache: an entry $(k,S_k)\mapsto\beta^{*}$ with $\beta'\le\beta^{*}$, where an
entry is stored only after $(k,S_k)$ was shown safe at $\beta^{*}$ by one of
the rules or by exhausting its children; rule~5 stores
$\beta^{*}=\mathrm{thr}(\pi(Q),hi)$.
\end{enumerate}
\end{theorem}
\begin{proof}
(1) The decided survivors are present in every completion. (2) Costs are
nonnegative, so a completion within $\beta'$ has $P_i(r_i)\le\beta'$ and, $P_i$
being nondecreasing, $r_i\le d_i$; at least $n_i-d_i$ copies of $v_i$ survive,
and a subset sum of a sub-multiset is a subset sum of the multiset. (3)--(5)
are Proposition~\ref{prop:window-rule}, Corollary~\ref{cor:root-rules} and
Theorem~\ref{thm:node-rule}. (6) is Lemma~\ref{lem:node-state} together with
Theorem~\ref{thm:node-rule}, which holds at every budget up to the stored
threshold.
\end{proof}
\noindent Not formalised as a whole: rules~(3)--(5) are the results cited in
the proof (rule~(3) only in part, see Proposition~\ref{prop:window-rule}), and
rules~(1), (2) and~(6) and Lemma~\ref{lem:node-state} are argued here.

\begin{theorem}[Exactness]\label{thm:kip-exact}
The search visits the root; at a node it returns ``safe'' if a rule of
Theorem~\ref{thm:kip-rules} applies or if all children
$r_{k+1}\in\{r\le n_{k+1}:P_{k+1}(r)\le\beta'\}$ return ``safe''; at a leaf
($k=L$) to which no rule applies it returns ``blocking'' with the counts on
the path. If the search ends without reaching its node limit, then ``safe''
at the root means the class instance is safe, and ``blocking'' returns a
blocking deletion.
\end{theorem}
\begin{proof}
By induction from the leaves, every node that returns ``safe'' is safe: for
the rules this is Theorem~\ref{thm:kip-rules}; for a node whose children are
all safe, every completion within $\beta'$ starts with some $r_{k+1}$ with
$P_{k+1}(r_{k+1})\le\beta'$ (costs are nonnegative), and the rest of it is a
completion of that child within $\beta'-P_{k+1}(r_{k+1})$. A leaf returning
``blocking'' has $S_L\cap[lo,hi]=\emptyset$ (rule~1 did not apply), and $M_L=R(r)$
for the path's deletion $r$, whose cost is $\beta-\beta'\le\beta$; so $r$ is
blocking. A node returns ``blocking'' only when a child does.
\end{proof}
\noindent Not formalised (the search code; see Section~\ref{app:lean-map}).
The rules it relies on are formalised as listed above, except the fractional
bound of Lemma~\ref{lem:fractional}.

\begin{corollary}[Exact solve]\label{cor:kip-loop}
Start from any leader set $X_0$ within budget with $u_0=\Phi(X_0)$ (the code
uses Weninger and Fukasawa's greedy leader). For $h=0,1,\dots$ run the search
with $\ell=u_h$: if it returns ``safe'', $\mathrm{OPT}=u_h$; if it returns a
blocking deletion $r$, set $X_{h+1}=X(r)$ and $u_{h+1}=\Phi(X(r))<u_h$. If no
search hits its node limit, this ends after at most $u_0+1$ searches with a
proof of optimality and an optimal leader set.
\end{corollary}
\begin{proof}
By Proposition~\ref{prop:kip-reduction}(3) and Theorem~\ref{thm:kip-exact},
``safe'' proves $\mathrm{OPT}\ge u_h=\Phi(X_h)\ge\mathrm{OPT}$, and ``blocking''
yields $X(r)$ within budget with $\Phi(X(r))<u_h$. (If
$\lceil u_h/g\rceil>hi$, no search is needed: $r=0$ is blocking, as noted
after Proposition~\ref{prop:kip-reduction}.) The values $u_h$ are
nonnegative integers that strictly decrease.
\end{proof}

For each of the \NumNewKip{} instances that were open in the literature the loop ended
with ``safe'' at BOBILib's own value, so BOBILib's point was already optimal
and the proof of optimality is new. These proofs are computational proofs
supported by formally verified lemmas, not formally verified solutions: Lean
proves the pruning rules (Section~\ref{app:lean-map}), and the search that
applies them, like the two integer-only provers that repeat it
(Section~\ref{app:verification}), is an exhaustive computation in integer
arithmetic whose code is tested but not verified. In the terms of
Section~\ref{app:classes} they are exact certificates (class~(i)).
